\documentclass[../../main.tex]{subfiles}

\begin{document}

\section{Bessel 不等式与规范正交基}\label{section:ch16-bessel-orthobasis}

在本节中, 设 $H$ 是内积空间.

\begin{definition}[正交集, 规范正交集]\label{definition:ch16-orthonormal}
$H$ 的子集 $S$ 称为正交的, 若 $S$ 中的每两个向量是正交的. 若这样的集合还具有性质: $S$ 中的每个向量是单位向量, 则 $S$ 称为规范正交的.
\end{definition}

\begin{theorem}[广义 Pythagorean 恒等式]\label{theorem:ch16-pythagorean}
若 $u_{1},\ u_{2},\ \cdots,\ u_{k}$ 是 $H$ 的规范正交向量, 而 $\alpha_{1},\ \cdots,\ \alpha_{n}$ 是实数, 则
\begin{align*}
\|\alpha_{1}u_{1}+\cdots+\alpha_{n}u_{n}\|^{2}=|\alpha_{1}|^{2}+\cdots+|\alpha_{n}|^{2}.
\end{align*}
\end{theorem}

该等式从展开以下等式
\begin{align*}
\|\alpha_{1}u_{1}+\cdots+\alpha_{n}u_{n}\|^{2}=\langle\alpha_{1}u_{1}+\cdots+\alpha_{n}u_{n},\alpha_{1}u_{1}+\cdots+\alpha_{n}u_{n}\rangle
\end{align*}
的右边得到.

\begin{theorem}[Bessel 不等式]\label{theorem:ch16-bessel}
对于 $H$ 中的规范正交序列 $\{\varphi_{k}\}$ 与 $H$ 中的向量 $h$,
\begin{align*}
\sum_{k=1}^{\infty}\langle\varphi_{k},h\rangle^{2}\leqslant\|h\|^{2}.
\end{align*}
\end{theorem}

为证明此不等式, 固定一个自然数 $n$ 且定义 $h_{n}=\sum_{k=1}^{n}\langle\varphi_{k},\ h\rangle\varphi_{k}$. 则根据广义的 Pythagorean 恒等式,
\begin{align*}
0\leqslant\|h-h_{n}\|^{2}=\|h\|^{2}-2\langle h,h_{n}\rangle+\|h_{n}\|^{2}=\|h\|^{2}-2\sum_{k=1}^{n}\langle h,\varphi_{k}\rangle\langle h,\varphi_{k}\rangle+\sum_{k=1}^{n}\langle h,\varphi_{k}\rangle^{2}=\|h\|^{2}-\sum_{k=1}^{n}\langle h,\varphi_{k}\rangle^{2}.
\end{align*}
因此
\begin{align*}
\sum_{k=1}^{n}\langle\varphi_{k},h\rangle^{2}\leqslant\|h\|^{2}.
\end{align*}
取 $n\to\infty$ 的极限得到 Bessel 不等式.

\begin{proposition}\label{proposition:ch16-fourier-series}
令 $\{\varphi_{k}\}$ 为 Hilbert 空间 $H$ 中的规范正交序列, 而向量 $h$ 属于 $H$. 则级数
\begin{align*}
\sum_{k=1}^{\infty}\langle\varphi_{k},\ h\rangle\varphi_{k}
\end{align*}
在 $H$ 中强收敛且向量 $h-\sum_{k=1}^{\infty}\langle\varphi_{k},\ h\rangle\varphi_{k}$ 正交于每个 $\varphi_{k}$.
\end{proposition}

\begin{proof}
对自然数 $n$, 定义 $h_{n}=\sum_{k=1}^{n}\langle\varphi_{k},\ h\rangle\varphi_{k}$. 根据广义的 Pythagorean 恒等式, 对每对自然数 $n$ 与 $k$,
\begin{align*}
\|h_{n+k}-h_{n}\|^{2}=\sum_{i=n+1}^{n+k}\langle\varphi_{i},h\rangle^{2}.
\end{align*}
然而, 根据 Bessel 不等式, 级数 $\sum_{k=1}^{\infty}\langle\varphi_{i},\ h\rangle^{2}$ 收敛, 因此 $\{h_{n}\}$ 是 $H$ 中的一个 Cauchy 列. 由于 $H$ 是完备的, $\sum_{k=1}^{\infty}\langle\varphi_{k},\ h\rangle\varphi_{k}$ 强收敛到向量 $h_{*}\in H$. 固定自然数 $m$. 观察到若 $n>m$, 则 $h-h_{n}$ 正交于 $\varphi_{m}$. 根据内积的连续性, $h-h_{*}$ 正交于 $\varphi_{m}$.
\end{proof}

\begin{definition}[完全规范正交序列]\label{definition:ch16-complete-ono}
Hilbert 空间 $H$ 中的规范正交序列 $\{\varphi_{k}\}$ 称为是完全的, 若 $H$ 中的向量仅有 $h=0$ 正交于每个 $\varphi_{k}$.
\end{definition}

我们从推论\cref{corollary:ch16-complete-span}推出 Hilbert 空间 $H$ 中的规范正交序列 $\{\varphi_{k}\}$ 是完全的当且仅当 $\{\varphi_{k}\}$ 的闭线性扩张是 $H$.

\begin{definition}[规范正交基]\label{definition:ch16-ono-basis}
Hilbert 空间 $H$ 中的规范正交序列 $\{\varphi_{k}\}$ 称为是 $H$ 的规范正交基, 若
\begin{align}
\text{对所有 }h\in H,\quad h=\sum_{k=1}^{\infty}\langle\varphi_{k},h\rangle\varphi_{k}.\label{align:ch16-fourier}
\end{align}
\end{definition}

\begin{proposition}\label{proposition:ch16-complete-iff-basis}
Hilbert 空间 $H$ 中的规范正交序列 $\{\varphi_{k}\}$ 是完全的当且仅当它是规范正交基.
\end{proposition}

\begin{proof}
首先假设 $\{\varphi_{k}\}$ 是完全的. 根据前一个命题, $h-\sum_{k=1}^{\infty}\langle\varphi_{k},\ h\rangle\varphi_{k}$ 正交于每个 $\varphi_{k}$. 因此, 根据 $\{\varphi_{k}\}$ 的完全性, \eqref{align:ch16-fourier}成立. 反之, 假定\eqref{align:ch16-fourier}成立. 若 $h\in H$ 正交于所有 $\varphi_{k}$, 则
\begin{align*}
h=\sum_{k=1}^{\infty}\langle\varphi_{k},h\rangle\varphi_{k}=\sum_{k=1}^{\infty}0\cdot\varphi_{k}=0.
\end{align*}
\end{proof}

\begin{example}\label{example:ch16-trig-system}
由取值为 $1/\sqrt{2\pi}$ 的常数函数与函数 $\{1/\sqrt{\pi}\cdot\sin kt,\ 1/\sqrt{\pi}\cdot\cos kt\}_{k=1}^{\infty}$ 组成的 $L^{2}[0,\ 2\pi]$ 中的函数的可数族是 Hilbert 空间 $L^{2}[0,\ 2\pi]$ 中的完全规范正交序列. 的确, 我们从初等的三角恒等式推出该序列是规范正交的. 我们从 Stone-Weierstrass 定理推出该序列的线性扩张关于最大值范数在 Banach 空间 $C[a,\ b]$ 中是稠密的. 因此根据 $C[a,\ b]$ 在 $L^{2}[0,\ 2\pi]$ 中的稠密性, 该序列的线性扩张在 $L^{2}[0,\ 2\pi]$ 中稠密.
\end{example}

若 Hilbert 空间 $H$ 具有规范正交基 $\{\varphi_{k}\}$, 则由于这些 $\varphi_{k}$ 的有限有理线性组合是 $H$ 的可数稠密子集, $H$ 必须是可分的. 事实证明可分性也是 Hilbert 空间具有规范正交基的充分条件.

\begin{theorem}\label{theorem:ch16-separable-ono-basis}
每个无穷维可分 Hilbert 空间具有规范正交基.
\end{theorem}

\begin{proof}
令 $\mathcal{F}$ 为 $H$ 的规范正交的子集族. 用包含关系将 $\mathcal{F}$ 排序. 对 $\mathcal{F}$ 的每个线性序子族, 该子族中集合的并是这个子族的一个上界. 根据 Zorn 引理, 我们可以选取 $\mathcal{F}$ 的最大子集 $S_{0}$. 由于 $H$ 是可分的, $S_{0}$ 是可数的. 令 $\{\varphi_{k}\}_{k=1}^{\infty}$ 为 $S_{0}$ 的列举. 若 $h\in H$, $h\neq 0$, 则根据命题\cref{proposition:ch16-fourier-series}, $h-\sum_{k=1}^{\infty}\langle\varphi_{k},\ h\rangle\varphi_{k}$ 正交于每个 $\varphi_{k}$. 因此 $h-\sum_{k=1}^{\infty}\langle\varphi_{k},\ h\rangle\varphi_{k}=0$, 否则 $S_{0}$ 与 $h-\sum_{k=1}^{\infty}\langle\varphi_{k},\ h\rangle\varphi_{k}$ 的规范化的并将是真包含 $S_{0}$ 的规范正交集. 因此 $\{\varphi_{k}\}_{k=1}^{\infty}$ 是 $H$ 的规范正交基.
\end{proof}

\begin{exercise}\label{exercise:ch16-19}
证明可分 Hilbert 空间 $H$ 的规范正交集必须是可数的.
\end{exercise}

\begin{exercise}\label{exercise:ch16-20}
令 $\{\varphi_{k}\}$ 为 Hilbert 空间 $H$ 的规范正交序列. 证明 $\{\varphi_{k}\}$ 在 $H$ 中弱收敛到 $0$.
\end{exercise}

\begin{exercise}\label{exercise:ch16-21}
令 $\{\varphi_{k}\}$ 为可分 Hilbert 空间 $H$ 的规范正交基. 证明在 $H$ 中 $\{u_{k}\}\to u$ 当且仅当 $\{u_{k}\}$ 是有界的且对每个 $k$, $\lim_{n\to\infty}\langle u_{n},\ \varphi_{k}\rangle=\langle u,\ \varphi_{k}\rangle$.
\end{exercise}

\begin{exercise}\label{exercise:ch16-22}
证明任何两个无穷维可分 Hilbert 空间是等距同构的且任何这样的同构保持内积.
\end{exercise}

\begin{exercise}\label{exercise:ch16-23}
令 $H$ 为 Hilbert 空间而 $V$ 是 $H$ 的闭可分子空间, 使得 $\{\varphi_{k}\}$ 是规范正交基. 证明 $H$ 映上 $V$ 的正交投影 $P$ 由
\begin{align*}
\text{对所有 }h\in H,\quad P(h)=\sum_{k=1}^{\infty}\langle\varphi_{k},h\rangle\varphi_{k}
\end{align*}
给出.
\end{exercise}

\begin{exercise}\label{exercise:ch16-24}
(Parseval 等式)令 $\{\varphi_{k}\}$ 为 Hilbert 空间 $H$ 的规范正交基. 证明:
\begin{align*}
\text{对所有 }h\in H,\quad\|h\|^{2}=\sum_{k=1}^{\infty}\langle\varphi_{k},h\rangle^{2}.
\end{align*}
也证明:
\begin{align*}
\text{对所有 }u,v\in H,\quad\langle u,v\rangle=\sum_{k=1}^{\infty}a_{k}\cdot b_{k},
\end{align*}
其中对每个自然数 $k$, $a_{k}=\langle u,\ \varphi_{k}\rangle$ 而 $b_{k}=\langle v,\ \varphi_{k}\rangle$.
\end{exercise}

\begin{exercise}\label{exercise:ch16-25}
证明关于 $L^{2}[0,\ 2\pi]$ 的规范正交基例子中的断言.
\end{exercise}

\begin{exercise}\label{exercise:ch16-26}
用命题\cref{proposition:ch16-complete-iff-basis}与 Stone-Weierstrass 定理证明: 对每个 $f\in L^{2}[-\pi,\ \pi]$,
\begin{align*}
f(x)=a_{0}/2+\sum_{k=1}^{\infty}[a_{k}\cdot\cos kx+b_{k}\cdot\sin kx],
\end{align*}
其中收敛依据 $L^{2}[-\pi,\ \pi]$ 范数且每个
\begin{align*}
a_{k}=\frac{1}{\pi}\int_{-\pi}^{\pi}f(x)\cos kx\,dx,\quad b_{k}=\frac{1}{\pi}\int_{-\pi}^{\pi}f(x)\sin kx\,dx.
\end{align*}
\end{exercise}

\end{document}